% ===========================================================
%  THEOREM 2 -- UNIFIED SHIFT DECOMPOSITION (proof)
%  v15 NeurIPS 2026 / ICML 2027 chapter
%  Author: Seungho Cook
% ===========================================================
\documentclass[11pt,a4paper]{article}
\usepackage[margin=1in]{geometry}
\usepackage{amsmath,amssymb,amsthm,bm,hyperref,booktabs,xcolor}
\newtheorem{theorem}{Theorem}
\newtheorem{proposition}{Proposition}
\newtheorem{corollary}{Corollary}
\newtheorem{lemma}{Lemma}
\newtheorem{definition}{Definition}
\newcommand{\R}{\mathbb{R}}
\newcommand{\E}{\mathbb{E}}
\newcommand{\KL}{\mathrm{KL}}
\newcommand{\TODO}[1]{\textbf{\textcolor{red}{[TODO: Seungho verify -- #1]}}}
\title{Theorem~2: Unified Shift Decomposition under Linear Correction\\
\large v15 NeurIPS/ICML chapter --- proof sketch (heavy TODO)}
\author{Seungho Cook}
\date{2026-04-25}
\begin{document}
\maketitle
\noindent\emph{This document must be read alongside
\texttt{theorem2\_setup.tex}.}

% ---------------------------------------------------------------
\section{Statement}

\begin{theorem}[Unified shift decomposition]
\label{thm:unified}
Assume class-conditionals are Gaussian,
$p_s(\tilde x\mid y) = \mathcal{N}(\tilde m_{s,y},\Sigma)$,
$p_t(\tilde x\mid y) = \mathcal{N}(\tilde m_{t,y},\Sigma)$ with
shared covariance $\Sigma\succ 0$. Let $T$ be a linear correction
operator (Definition~\ref{def:linear-correction} of the setup
document) and let $w_Y$ be the population Fisher direction under
$p_s$. Let $P_Y = w_Y w_Y^{\top}/\|w_Y\|^{2}$ denote the rank-one
projector onto the label axis and $P_Y^{\perp}=I-P_Y$ its
orthogonal complement. Then:
\begin{equation}
  \mathrm{DIAL}(\tilde X,Y)
  \;=\; \Phi\!\bigl(\Delta_{\text{cov}},\,\Delta_{\text{lab}},\,
                     \Delta_{\text{cond}}^{\parallel},\,
                     \Delta_{\text{cond}}^{\perp}\bigr),
  \label{eq:thm2}
\end{equation}
where
\[
   \Delta_{\text{cond}}^{\parallel}
   = \E_{y}\!\left[\tfrac12\bigl\|P_Y\,\Sigma^{-1/2}
     (\tilde m_{s,y}-\tilde m_{t,y})\bigr\|^{2}\right],\qquad
   \Delta_{\text{cond}}^{\perp}
   = \E_{y}\!\left[\tfrac12\bigl\|P_Y^{\perp}\,\Sigma^{-1/2}
     (\tilde m_{s,y}-\tilde m_{t,y})\bigr\|^{2}\right],
\]
and $\Phi$ is \emph{strictly increasing in $\Delta_{\text{cond}}^{\parallel}$
and independent of $\Delta_{\text{cov}}$ and $\Delta_{\text{lab}}$} (up
to a $1/\sqrt{n}$ CV noise term).
\end{theorem}

\begin{corollary}[Subspace alignment]
\label{cor:subspace}
DIAL $>$ 0 if and only if
$\Delta_{\text{cond}}^{\parallel} > \Delta_{\text{cond}}^{\perp}$, i.e.\
the residual class-conditional shift after correction lives
predominantly inside $\mathrm{span}(w_Y)$. Equivalently, the flip
condition of Theorem~1 (v10 chapter) is the event $\alpha>1/2$ when
the batch contamination is itself rank-one; Theorem~2 generalises it
to arbitrary rank.
\end{corollary}

\begin{corollary}[Unification of shift taxonomy]
\label{cor:unify}
Let $T$ be a linear correction producing $\tilde X=T(X,B)$. Then for
the canonical shift types:
\begin{enumerate}
\item[\textup{(a)}] \emph{Pure covariate shift}
  ($p_s(x)\ne p_t(x),$ same $p(y\mid x)$):
  $\Delta_{\text{cov}}>0,\ \Delta_{\text{cond}}^{\parallel}\to 0$
  as $T$ converges to the Bayes-optimal importance-weighting map, so
  $\mathrm{DIAL}\to 0$.
\item[\textup{(b)}] \emph{Pure label (prior) shift}: $\Delta_{\text{lab}}>0$
  only, $\Delta_{\text{cond}}^{\parallel}=0$, thus
  $\mathrm{DIAL}=0$.
\item[\textup{(c)}] \emph{Subspace-aligned conditional shift}:
  $\Delta_{\text{cond}}^{\parallel}>0$ and
  $\mathrm{DIAL}>0$. This is the regime uniquely diagnosed by our
  metric.
\end{enumerate}
\end{corollary}

% ---------------------------------------------------------------
\section{Proof sketch}

We prove \eqref{eq:thm2} in six steps.

% ------------------------------------
\paragraph{Step 1 (linear correction as affine projection).}
Let $T(x,b)=A_b x + c_b$. For each batch $b$, define the \emph{residual
batch subspace}
$\mathcal{S}_b = A_b\,\mathcal{X}\,\subset\R^{p}$.
Because $A_b$ is non-singular, $\mathcal{S}_b=\R^{p}$ as a set, but the
induced metric is $A_b\,\Sigma A_b^{\top}$; up to change of basis we
may assume, WLOG, that $A_b=I$ and $c_b$ carries all the correction
(mean-centring case). This absorbs the scale-component of ComBat into
$\Sigma$. \TODO{The reduction $A_b\!=\!I$ is WLOG only if $A_b$ commutes
with the Bayes classifier's whitening matrix; verify this is valid for
ComBat's diagonal $A_b$, and that the general case merely rescales
DIAL by a constant $\kappa(A)\in(0,1]$ (state as a lemma in the final
version).}

% ------------------------------------
\paragraph{Step 2 (KL chain rule decomposition).}
Write the joint under each distribution as
$p(\tilde x,y)=p(y)\,p(\tilde x\mid y)$. The KL divergence between
the \emph{joint corrected} distributions factorises as:
\begin{align*}
  \KL\!\bigl(p_s(\tilde x,y)\,\|\,p_t(\tilde x,y)\bigr)
   &= \KL\!\bigl(p_s(y)\,\|\,p_t(y)\bigr)
      + \E_{y\sim p_s}\bigl[
        \KL\!\bigl(p_s(\tilde x\mid y)\,\|\,
                   p_t(\tilde x\mid y)\bigr)\bigr]\\
   &= \Delta_{\text{lab}} + \Delta_{\text{cond}}.
\end{align*}
By the marginal--conditional identity
$p(\tilde x) = \int p(\tilde x\mid y)p(y)\,dy$ and Csisz\'ar's data
processing inequality,
$
  \Delta_{\text{cov}} \;\le\; \Delta_{\text{lab}} + \Delta_{\text{cond}}.
$
Equality holds iff $p_s(y\mid \tilde x)=p_t(y\mid \tilde x)$. \TODO{Write
the chain rule decomposition completely: need to justify the swap of
integration and KL for Gaussian mixtures, and cite Cover \& Thomas
Thm.~2.5.3.}

% ------------------------------------
\paragraph{Step 3 (class-conditional KL under shared covariance).}
For Gaussians with shared $\Sigma$,
\[
  \KL\!\bigl(\mathcal{N}(\tilde m_{s,y},\Sigma)\,\|\,
             \mathcal{N}(\tilde m_{t,y},\Sigma)\bigr)
  \;=\; \tfrac12\,
        (\tilde m_{s,y}-\tilde m_{t,y})^{\!\top}
        \Sigma^{-1}
        (\tilde m_{s,y}-\tilde m_{t,y}).
\]
Introducing the projectors $P_Y,P_Y^{\perp}$ on the pre-whitened
direction $\Sigma^{-1/2}w_Y$,
\[
  \Delta_{\text{cond}}
  \;=\; \underbrace{\E_y[\tfrac12\|P_Y\,\Sigma^{-1/2}
        (\tilde m_{s,y}-\tilde m_{t,y})\|^{2}]}_{=\Delta_{\text{cond}}^{\parallel}}
      + \underbrace{\E_y[\tfrac12\|P_Y^{\perp}\,\Sigma^{-1/2}
        (\tilde m_{s,y}-\tilde m_{t,y})\|^{2}]}_{=\Delta_{\text{cond}}^{\perp}}.
\]

% ------------------------------------
\paragraph{Step 4 (Bayes-classifier posterior under corrected means).}
For Gaussian classes with shared $\Sigma$ the log-odds
$\log p_s(y{=}1\mid\tilde x)/p_s(y{=}0\mid\tilde x)$ is the linear
form $\langle w_Y,\tilde x\rangle + c$ up to constants, and the
\emph{flipped} log-odds on $p_t$ is
$\langle w_Y',\tilde x\rangle + c'$ with
$w_Y' = \Sigma^{-1}(\tilde m_{t,1}-\tilde m_{t,0})$. Hence
$\mathrm{sign}\,\mathrm{AUC}(\tilde X_t,Y_t)-\tfrac12
 = \mathrm{sign}\,\langle w_Y,w_Y'\rangle$. When the shift that has
survived correction lies in $\mathrm{span}(w_Y)$, $w_Y'\approx-w_Y$
and AUC $<\tfrac12$, giving DIAL $>$ 0. The magnitude of the flip is
monotone in $\Delta_{\text{cond}}^{\parallel}$. \TODO{Formalise the
monotonicity: compute $\partial\mathrm{AUC}/\partial
\Delta_{\text{cond}}^{\parallel}$ explicitly in the Gaussian equal-covariance
case, using Das~Gupta (2008) Prop.~4.3.}

% ------------------------------------
\paragraph{Step 5 (CV noise term).}
$\mathrm{AUC}$ is evaluated by $k$-fold cross-validation so observed
DIAL equals the population-DIAL plus a mean-zero noise term of
magnitude $O_p(\sqrt{\log(1/\delta)/n})$ by McDiarmid's inequality
applied to the U-statistic form of AUC (Agarwal \emph{et al.}\ 2005,
Thm.~3). This contributes the additive $\varepsilon_n$ in
\eqref{eq:thm2}. \TODO{Cite McDiarmid with the exact bound constant;
current expression is informal.}

% ------------------------------------
\paragraph{Step 6 (assembly).}
Combining Steps 2--5 gives Theorem~\ref{thm:unified} with
\[
   \Phi(\cdot) \;\propto\;
   \mathbf{1}\!\bigl[\Delta_{\text{cond}}^{\parallel}
                     >\Delta_{\text{cond}}^{\perp}\bigr]
   \cdot
   \min\!\bigl(1,\,c_0\,\Delta_{\text{cond}}^{\parallel}\bigr),
\]
for an absolute constant $c_0>0$ depending on $\|\Sigma^{-1/2}w_Y\|$.
The claim that $\Phi$ is \emph{independent of
$\Delta_{\text{cov}}$ and $\Delta_{\text{lab}}$ up to $\varepsilon_n$}
follows because AUC is invariant to strictly monotone rescalings of
the score --- a covariate-only shift or prior-only shift rescales both
class scores identically, leaving the classifier's ranking unchanged.
\TODO{This AUC-invariance step is the crux of the unification. Double
check the rank-preservation argument with a counter-example run in
\texttt{v15\_theorem2\_verify.py} rows labelled ``shift\_type=label''.}

\qed

% ---------------------------------------------------------------
\section{Connections}

\paragraph{Ben-David \emph{et al.}\ (2010) $\mathcal{H}$-divergence.}
The A-distance $d_{\mathcal{A}}(p_s,p_t)$ upper-bounds the
transferability risk by
$\epsilon_t\le\epsilon_s+d_{\mathcal{H}\Delta\mathcal{H}}+\lambda$.
Theorem~2 shows that when $T$ is linear and alignment $\alpha>1/2$,
the A-distance is dominated by the parallel component
$\Delta_{\text{cond}}^{\parallel}$, giving DIAL as a computable
surrogate for the otherwise intractable
$\mathcal{H}\Delta\mathcal{H}$-divergence. \TODO{Exact inequality:
$\Delta_{\text{cond}}^{\parallel}\le d_{\mathcal{H}\Delta\mathcal{H}}
\le 2\sqrt{\Delta_{\text{cond}}}$ --- prove in appendix.}

\paragraph{Sugiyama \emph{et al.}\ (2007) KLIEP.}
KLIEP estimates the importance weight
$r(x)=p_t(x)/p_s(x)$ by minimising an empirical KL. Theorem~2 says
that KLIEP cannot correct DIAL $>$ 0 alone: when
$\Delta_{\text{cond}}^{\parallel}>0$, re-weighting the source cannot
restore the posterior direction, since re-weighting acts on
$p(x)$ not $p(x\mid y)$.

\paragraph{Arjovsky \emph{et al.}\ (2019) IRM.}
IRM searches for representations $\Phi(X)$ under which the classifier
is simultaneously optimal across environments. Theorem~2 implies that
IRM's \emph{failure mode} (partial invariance) is precisely the regime
where $\Delta_{\text{cond}}^{\parallel}>0$ --- i.e.\ DIAL is positive
if and only if IRM has not reached an invariant solution. This gives
IRM a new empirical diagnostic.

\paragraph{Tishby \emph{et al.}\ (2015) information bottleneck.}
Let $I(\tilde X;Y)$ be the mutual information of corrected features
with labels. The flip event DIAL $>$ 0 is equivalent (up to
$\varepsilon_n$) to $I(\tilde X;Y\mid B)\ll I(X;Y)$, i.e.\ correction
has \emph{destroyed information} about $Y$ while preserving information
about $B^\perp$. \TODO{Make precise: show the identity
$I(\tilde X;Y\mid B) = I(X;Y) - \Delta_{\text{cond}}^{\parallel}+O(1/n)$
under our Gaussian assumptions.}

\section{Numerical validation}

See companion script
\texttt{\detokenize{notebooks_or_scripts/v15_theorem2_verify.py}},
which simulates
\[
   \{\Delta_{\text{cov}},\,\Delta_{\text{lab}},\,
    \Delta_{\text{cond}}^{\parallel},\,
    \Delta_{\text{cond}}^{\perp}\}
\]
independently and regresses observed DIAL against the decomposition.
Results file: \texttt{results/v15\_neurips/theory\_validation/
theorem2\_numerical.tsv}; interactive visualisation
\texttt{reports/html/figs\_interactive/v15/
theorem2\_decomposition.html}.

\paragraph{Expected result.}
Under the four ablation regimes (pure covariate, pure label, rank-one
subspace-aligned, full conditional) the sign and magnitude of DIAL
should match the prediction of \eqref{eq:thm2} with
$R^{2}\gtrsim 0.9$. If this $R^{2}$ is materially lower, one of the
TODO steps above is the likely culprit.

\end{document}
